% FORMALIZACION_REUNIDA: composición de la respuesta constitutiva de III
% con el lector positivo de GEOMETRIA_POSITIVA/08_ARTICULO_AUTONOMO.tex.
% Propietario de las operaciones: III/manuscrito__sections__04_respuesta_constitutiva.tex.
% Residencia propuesta: después de la coherencia multimorfológica, antes del
% empuje amplituhedral. No modifica sus hipótesis de frontera, residuo o grado.
\subsection{Realización constitutiva positiva y naturalidad espectral}
\label{empalme:positividad-constitutiva}

La realización constitutiva de la
proposición~\ref{geometria_positiva:prop:multimorphology-coherence} recibe
un operador ya construido desde las dos hojas orientadas del TPK. La
sección~\ref{iii:iii:sec:constitutiva} determina su respuesta y su inversa;
la presente composición explicita la forma positiva y el enlace de
refinamiento que utiliza esa realización.

Sea $\widehat{\mathcal X}_{\rm con}$ el dominio de estados enriquecidos
cuya lectura angular pertenece a la cámara $x>y>0$. Para
$\xi\in\widehat{\mathcal X}_{\rm con}$, la lectura produce un espacio
de dos hojas $E_\xi$, una involución ortogonal $\mathcal R_\xi$ y el
transporte
\[
 P_{\pm,\xi}=\frac{I\pm\mathcal R_\xi}{2},\qquad
 T_\xi=q_{+,\xi}P_{+,\xi}+q_{-,\xi}P_{-,\xi},
 \qquad 0<q_{+,\xi}<q_{-,\xi}<1.
\]
Las incidencias $90$ y $120$ proceden del calendario y actúan sobre
ese mismo transporte. La respuesta es
\[
 V_\xi=(I-T_\xi^{90})(I-T_\xi^{120})^{-1},\qquad
 W_{m,\xi}=V_\xi\otimes I_{9^m}.
\]
La amplificación conserva todas las posiciones de la nueva fibra
nonádica. Sobre $E_{m,\xi}=E_\xi\otimes\mathbb C^{9^m}$ se define
la forma
\begin{equation}
 \mathfrak C_m(\xi)(v,w)
 :=\langle W_{m,\xi}v,W_{m,\xi}w\rangle
 =\langle v,W_{m,\xi}^{2}w\rangle.
 \label{empalme:eq:forma-constitutiva}
\end{equation}
Ésta es una realización concreta en el cono de formas hermíticas
positivas definidas. El cuadrado conserva las respuestas cuadráticas
de las dos hojas que publican las coordenadas
$\widehat\varepsilon$ y $\widehat\mu$ de
\eqref{iii:iii:eq:four}.

\begin{proposition}[Coercividad, recuperación y enlace nonádico]
\label{empalme:prop:constitutiva-natural}
La forma~\eqref{empalme:eq:forma-constitutiva} satisface, para todo
$v\in E_{m,\xi}$ no nulo,
\[
 \frac9{16}\|v\|^2
 <\mathfrak C_m(\xi)(v,v)<\|v\|^2.
\]
Si $A_{m,\xi}=W_{m,\xi}^{2}$ y $\mathbb E_m$ es la expectativa
parcial normalizada sobre el último factor de dimensión nueve, entonces
\[
 \mathbb E_m(A_{m+1,\xi})=A_{m,\xi}.
\]
La forma, junto con las etiquetas orientadas de hoja, recupera
$V_\xi$, los dos canales $q_{\pm,\xi}$ y sus coordenadas angulares.
Si un transporte unitario $J:E_\xi\to E_{\xi'}$ satisface
$JT_\xi=T_{\xi'}J$, entonces
\[
 (J\otimes I)^*A_{m,\xi'}(J\otimes I)=A_{m,\xi}.
\]
La identidad se mantiene para cualquier función continua del espectro
que se aplique a la respuesta antes de amplificar.
\end{proposition}

\begin{proof}
Los proyectores de hoja son ortogonales y suman la identidad. El cálculo
espectral de~\eqref{iii:iii:eq:rchannels} da
\[
 V_\xi=r_{+,\xi}P_{+,\xi}+r_{-,\xi}P_{-,\xi},\qquad
 r_{\pm,\xi}=f(q_{\pm,\xi}^{30}),\qquad
 f(s)=\frac{1+s+s^2}{1+s+s^2+s^3}.
\]
En $0<s<1$ se tiene
\[
 4(1+s+s^2)-3(1+s+s^2+s^3)
 =(1-s)(3s^2+2s+1)>0,
\]
y el denominador supera al numerador en $s^3>0$. Por tanto
$3/4<r_{\pm,\xi}<1$. La descomposición ortogonal y el cuadrado de
estos límites prueban la coercividad, que permanece idéntica bajo
amplificación.

Como $A_{m+1,\xi}=A_{m,\xi}\otimes I_9$ y la traza normalizada
del último factor satisface $\operatorname{tr}_9(I_9)=1$,
\[
 \mathbb E_m(A_{m+1,\xi})
 =(\operatorname{id}\otimes\operatorname{tr}_9)
 (A_{m,\xi}\otimes I_9)=A_{m,\xi}.
\]
La misma igualdad vale para cada polinomio en $W_{m,\xi}$. La
aproximación uniforme sobre el espectro compacto y la continuidad de
la expectativa la extienden a toda función continua.

La raíz positiva de $A_{0,\xi}$ recupera $V_\xi$. Los proyectores
orientados identifican sus dos autovalores. La función $f$ es
estrictamente decreciente en $(0,1)$ por
\eqref{iii:iii:eq:monotone}; su inversa determina las raíces únicas
de~\eqref{iii:iii:eq:cubic}. La raíz positiva de orden treinta
restituye $q_{\pm,\xi}$, y
\eqref{iii:iii:eq:recover-angular} restituye $x$ e $y$.

Finalmente, $JT_\xi=T_{\xi'}J$ implica la misma identidad para
todas sus potencias. La invertibilidad de $I-T^{120}$ proporciona
\[
 J(I-T_\xi^{120})^{-1}
 =(I-T_{\xi'}^{120})^{-1}J.
\]
De aquí resulta $JV_\xi=V_{\xi'}J$. La amplificación, el cuadrado
y la unitariedad de $J$ dan la congruencia declarada. El argumento
polinómico y su extensión uniforme prueban la última afirmación.
\end{proof}

El enlace de esta construcción es la expectativa normalizada que
conserva el operador amplificado. La compatibilidad con un transporte
geométrico utiliza el entrelazamiento $JT_\xi=T_{\xi'}J$ expresado
en el enunciado. Por su parte, el transporte de formas canónicas y
residuos mantiene los cuadrados de frontera de la
proposición~\ref{geometria_positiva:prop:multimorphology-coherence}.
Así quedan precisadas las operaciones de cada nivel: respuesta angular,
forma constitutiva positiva, amplificación de la fibra y transporte de
incidencias. La restitución angular recupera esta lectura; el archivo
enriquecido conserva además la historia que la produjo.
